\ncdlesson{3}{Urbaniseringens pris: Hvem har råd til at bo i byen?}{城市化的代价：谁住得起大城市？}
{DU3 Modul 5（PD3）$\rightarrow$ Modul 6（Studieprøven, B2 $\rightarrow$ C1）｜对标 NCE 4 \textit{Samfundsanalyse \& Boligpolitik}}
{\textbf{悖论与动态同步结构}：反讽/悖论导入（\textit{paradoksalt nok}）· 动态同步（\textit{i takt med at...}）· 正反馈螺旋（\textit{en selvforstærkende effekt}）· 比例递进（\textit{jo mere ..., desto ...}）· 代价取舍（\textit{på bekostning af}）}
{在谈论哥本哈根或奥胡斯的高房价时，B1 水平的表达往往停留在“房子太贵、工资不够”。而要写出打动 Studieprøven 考官的社会学深度，你必须学会描绘事物之间的\textbf{动态联动与制度悖论}：城市越宜居、越绿色，反而越把普通劳动者挤出市中心；通勤距离越远，社会分化就越剧烈。本课将带你掌握 \textbf{\textit{paradoksalt nok}}、\textbf{\textit{i takt med at}} 与 \textbf{\textit{jo ..., desto ...}} 的严密句法，精准刻画城市化背后的自我强化循环。}

% ---------------------------------------------------------------------
\begin{foerduleaeser}
带着下面两个宏观问题阅读 Tekst A 和 Tekst B：
\begin{enumerate}
  \item \textit{Hvilket paradoks opstår der ifølge Tekst A, når storbyer som København bliver grønnere og mere attraktive, og hvordan skaber boligmarkedet en selvforstærkende effekt?}（根据 Tekst A，当哥本哈根等大城市变得更加绿色宜居时，会产生什么悖论？房地产市场又是如何制造“自我强化效应”的？）
  \item \textit{Hvordan adskiller Henriks og Sofies forslag til at løse boligpresset sig fra hinanden i byrådsdebatten i Tekst B?}（在 Tekst B 的市政议会辩论中，亨里克和索菲针对缓解住房压力的方案有何分歧，又在何处达成妥协？）
\end{enumerate}
\end{foerduleaeser}

% ---------------------------------------------------------------------
\ncdsubsec{二、Tekster（城市社会学分析与市政对谈双语篇）}

\subsection*{Tekst A：\textit{Når attraktive byer prissætter deres egne borgere ud}（城市社会学分析专栏 · 约 345 词）}

\begin{tipbox}{语篇说明}
本篇剖析丹麦大城市人口集聚背后的住房危机与通勤不平等。阅读时请特别关注\textbf{\color{NCDNavy}深蓝色标记的动态同步与比例递进句型}（\textit{i takt med at..., jo mere..., desto...}），以及\textbf{\color{NCDRed}深红色标记的悖论、反馈循环与代价取舍表达}（\textit{paradoksalt nok, en selvforstærkende effekt, på bekostning af}）。
\end{tipbox}

I mere end to årtier har Danmark oplevet en markant urbanisering, hvor både studerende, børnefamilier og vidensvirksomheder søger mod København og Aarhus. Ved første øjekast ligner udviklingen en ubetinget succeshistorie med grønne cykelstier, rene havnebade, moderne metroforbindelser og et blomstrende kulturliv. Men \da{i takt med at} storbyerne er blevet stadig mere attraktive og klimavenlige, er boligpriserne stukket af fra helt almindelige lønninger, så adgangen til byen nu afhænger af formue frem for behov.

Her opstår der et slående bypolitisk dilemma. \dabold{Paradoksalt nok} er det netop de vellykkede investeringer i byrum, kollektiv trafik og grønt miljø, der nu truer byens sociale mangfoldighed. Når nedslidte industrikvarterer omdannes til eftertragtede boligområder langs havnen, stiger grundværdierne eksplosivt, og sygeplejersker, pædagoger samt politibetjente har simpelthen ikke længere råd til at bo tæt på de borgere og arbejdspladser, de servicerer i hverdagen.

På boligmarkedet sætter denne udvikling hurtigt \dabold{en selvforstærkende effekt} i gang. \da{Jo mere} eksklusivt et kvarter fremstår, \da{desto} større bliver interessen fra velstillede købere og private pensionskasser, hvilket igen presser kvadratmeterpriserne endnu højere op. Samtidig prioriterer kommuner og udviklere ofte dyre ejerboliger eller luksuriøse lejeboliger, som giver hurtige indtægter og stærke skatteydere, \dabold{på bekostning af} billige almene boliger med en husleje til at betale.

Konsekvensen mærkes langt uden for kommunegrænsen i form af voksende pendlerulighed. Unge familier og ansatte i servicefag fortrænges gradvist til oplandskommuner, hvor boligerne stadig er til at betale. \da{Jo længere} væk fra arbejdspladsen disse borgere presses, \da{desto} flere timer og kroner sluger den daglige transport gennem morgenens trafikpropper. Dermed opstår der ikke blot voksende økonomiske forskelle, men også en skæv fordeling af fritid, energi og familietid mellem centrum og opland.

Hvis storbyerne ikke skal ende som lukkede enklaver for de mest velhavende, kræves der en langt mere målrettet bolig- og planpolitik. En aktiv kommunal jordpolitik, faste kvoter for almene boliger ved nybyggeri samt hurtige togforbindelser til omegnen kan bremse segregeringen og sikre blandede kvarterer. Kun ved at tænke grøn vækst og social balance sammen kan byen bevare plads til alle de mennesker, der får dens skoler, hospitaler og gader til at fungere.

\begin{transbox}
\small
\textbf{\color{NCDTeal}[Tekst A 中文精译]}\\
二十多年来，丹麦经历了显著的城市化进程，大学生、有孩子的家庭以及知识密集型企业纷纷向哥本哈根和奥胡斯集聚。乍看之下，这一发展宛如一个毫无瑕疵的成功故事：绿意盎然的自行车道、洁净的港湾泳池、现代化的地铁线路以及繁荣的文化生活。然而，随着大城市变得愈发具有吸引力且更加气候友好，房价却早已远远甩开了普通人的薪资水平，以至于能否进入城市生活如今更多取决于资产而非实际需求。

一个引人注目的城市政策困境由此浮现。令人深感悖论的是，恰恰是那些对城市空间、公共交通和绿色生态环境的成功投资，如今正在威胁着城市原本的社会多样性。当破旧的老工业区沿着港湾被改造成备受追捧的住宅区时，土地价值呈爆发式飙升，护士、幼教老师以及警察等群体干脆再也负担不起住在自己日常所服务的市民与工作岗位附近。

在房地产市场上，这种演变迅速启动了一种自我强化效应。一个街区显得越高端稀缺，富裕买家和私营养老基金的兴趣就越浓厚，而这反过来又将每平方米房价推得更高。与此同时，市镇政府和开发商往往优先建设能带来快速收益和高收入纳税人的高价产权房或奢华租赁公寓，却以牺牲租金可负担的平价社会福利住房为代价。

这一后果以日益扩大的通勤不平等形式，在市镇边界之外的广大地区被深切感受到。年轻家庭和服务业从业者逐渐被挤向房价尚可负担的周边卫星市镇。这些市民被推离工作地点越远，穿过早高峰拥堵车流的日常通勤吞噬的时间和金钱就越多。由此产生的不仅是不断扩大的经济差距，还有市中心与外围地区之间在闲暇时间、精力与家庭陪伴上的失衡分配。

如果不想让大城市最终沦为只服务于最富裕阶层的封闭飞地，就必须推行更有针对性的住房与城市规划政策。积极的市政土地政策、新建项目中固定的社会福利住房配额，以及通往周边地区的快捷轨道交通，都能够遏制社会隔离并保障混合街区。唯有将绿色增长与社会平衡统筹考虑，城市才能为所有维系其学校、医院和街道运转的人留出容身之地。
\end{transbox}

\subsection*{Tekst B：\textit{Debat fra byrådssalen: Vækst og blandede bydele}（市政议会政策对谈 · 约 155 词）}

\textbf{阅读任务（理解与核对）：}\textit{Læs debatten mellem finansordfører Henrik og byplanlægger Sofie, og besvar spørgsmålene nedenfor.}

\begin{emailbox}{市政辩论对谈｜Henrik (Finansordfører) \& Sofie (Byplanlægger)}
\textbf{Henrik:} Vi ønsker jo alle en levende by, men jeg er bekymret for forslaget om at tvinge udviklere til at bygge 30 procent almene boliger i alle nye projekter. \da{Jo strammere} krav kommunen stiller til private investorer, \da{desto} færre boliger bliver der i sidste ende bygget --- og så stiger priserne bare endnu mere på grund af boligmangel.

\textbf{Sofie:} Den risiko skal vi selvfølgelig tage alvorligt, Henrik. Men \dabold{paradoksalt nok} hjælper et stort nybyggeri af dyre lejligheder ikke en sygeplejerske, hvis huslejen ligger over 16.000 kroner. Når vi bygger udelukkende efter markedets højeste afkast, sker det \dabold{på bekostning af} byens sammenhængskraft.

\textbf{Henrik:} Det anerkender jeg delvist. Måske ligger kompromiset i en hurtigere kommunal sagsbehandling, så vi udbygger både private og almene projekter i takt med befolkningstilvæksten --- især omkring nye S-togsstationer i forstæderne.

\textbf{Sofie:} Et rigtig godt bud! Kombinerer vi blandede ejerformer i centrum med bedre regional transport, kan vi nemlig bryde den \dabold{selvforstærkende effekt} uden at bremse byens udvikling.
\end{emailbox}

\textbf{Tekst B 阅读核对（Rigtigt / Forkert / Kort svar，答案见第六部分）：}
\begin{enumerate}
  \item \textit{Henrik frygter, at strenge krav om 30 procent almene boliger vil føre til færre nybyggede boliger i alt.} \underline{\hspace{2.2cm}}
  \item \textit{Hvorfor mener Sofie ikke, at mange nye, dyre lejligheder hjælper almindelige lønmodtagere som sygeplejersker?} \underline{\hspace{2.2cm}}
  \item \textit{Hvilket konkret kompromis foreslår Henrik for at sikre både flere boliger og hurtigere udbygning?} \underline{\hspace{2.2cm}}
\end{enumerate}

\begin{tipbox}{本课发音与语调要点（Udtale \& Prosodi）}
\begin{itemize}
  \item \textbf{城市化与住房政策词汇}：\textit{urbanisering} \ipa{uɐ̯bɑniˈseːˀɐeŋ}（重音在 \textit{-se-} 且带喉塞音，\textit{-ering} 读 \ipa{ˈeːˀɐeŋ}）；\textit{en kvadratmeterpris} \ipa{kvɑˈd̥ʁɑːd̥ˌmeːˀdɐˌpʰʁiːˀs}；\textit{almen} \ipa{ˈalˌmeːˀn}（常见复数 \textit{almene boliger} \ipa{ˈalˌmeːˀnə ˈb̥oːliɐ}，注意 \textit{bolig} 中 \textit{g} 不发音！）。
  \item \textbf{动态结构词的静音字母与 \textit{stød}}：\textit{paradoksalt} \ipa{pʰɑʁɑd̥ʌɡˈsæːˀld}（重音在最末音节）；\textit{selvforstærkende} \ipa{ˈsɛlˀfʌˌsd̥ɛɐ̯ɡənə}（\textit{selv} 末尾 \textit{v} 常省略不读！）；\textit{på bekostning af} \ipa{pʰɔ beˈkʰʌsd̥neŋ ˈɑʊ̯ˀ}（\textit{af} 读 \ipa{ɑʊ̯ˀ}）。
  \item \textbf{比例比较句语调起伏}：朗读 \textit{Jo mere X, desto mere Y} 时，前半句在 \textit{Jo mere} 处微扬，到逗号前略作悬停；后半句在 \textit{desto} 后的比较级重音陡然落地，凸显因果递增张力。
\end{itemize}
\end{tipbox}

\begin{transbox}
\small
\textbf{\color{NCDTeal}[Tekst B 中文对照]}\\
\textbf{亨里克（财政发言人）：}我们当然都希望拥有一座充满活力的城市，但我对在所有新项目中强制要求开发商建设 30\% 社会福利住房的提案感到担忧。市政府对私营投资者施加的要求越严苛，最终建成的住房总数就会越少——而这样一来，房价只会因为房源短缺而涨得更凶。

\textbf{索菲（城市规划师）：}我们当然应当认真对待这种风险，亨里克。但具有讽刺意味的悖论在于，如果月租超过一万六千克朗，哪怕新建再多高价公寓也帮不到一名护士。当我们完全追逐市场最高回报率来建房时，其代价往往就是城市的社会凝聚力。

\textbf{亨里克：}这一点我部分认可。也许折中方案在于加快市政审批流程，使我们能够随着人口增长同步推进私营项目与社会福利住房项目——尤其是在郊区的新城铁车站周边。

\textbf{索菲：}这是一个非常好的建议！只要我们把市中心的多元产权混合模式与更完善的区域公共交通结合起来，就能在不扼杀城市发展的前提下打破这种自我强化的恶性循环。
\end{transbox}

% ---------------------------------------------------------------------
\ncdsubsec{三、Gloser（核心生词表）}

\begin{tipbox}{词表重点}
丹麦社会特有的 \da{almen bolig}（复数常说 \textit{almene boliger}）并非中文语境下简陋的“救济房”，而是由非营利住房协会（\textit{boligselskaber}）运营、面向广大普通工薪阶层的平价福利住房，是理解丹麦城市政策辩论的关键概念！此外，动词 \da{at prissætte}（定价）与英文短语 \textit{price out} 对应为 \textit{at prissætte nogen ud}（因高价将某群体挤出）。
\end{tipbox}

\begin{glosertable}
名词 & en urbanisering & \textit{(urbaniseringen, kun ental)} & 城市化，都市化 \\
名词 & et boligmarked & \textit{(boligmarkedet, boligmarkeder, boligmarkederne)} & 房地产市场，住房市场 \\
名词 & en kvadratmeterpris & \textit{(kvadratmeterprisen, kvadratmeterpriser, kvadratmeterpriserne)} & 每平方米房价 \\
名词 & en mangfoldighed & \textit{(mangfoldigheden, kun ental)} & 多样性，多元共存 \\
名词 & en byfornyelse & \textit{(byfornyelsen, byfornyelser, byfornyelserne)} & 城市更新，旧区改造 \\
名词 & en oplandskommune & \textit{(oplandskommunen, oplandskommuner, oplandskommunerne)} & 周边市镇，外围卫星城 \\
名词 & en segregering & \textit{(segregeringen, segregeringer, segregeringerne)} & 社会隔离，阶层分区居住 \\
名词 & en ejerform & \textit{(ejerformen, ejerformer, ejerformerne)} & 住房产权形式 \\
名词 & en sagsbehandling & \textit{(sagsbehandlingen, sagsbehandlinger, sagsbehandlingerne)} & 行政审批，案件处理 \\
名词 & en effekt & \textit{(effekten, effekter, effekterne)} & 效应，效果 \\
\midrule
动词 [har] & at prissætte & \textit{(prissætter, prissatte, har prissat)} & 给……定价（\textit{ud} 挤出） \\
动词 [har] & at fortrænge & \textit{(fortrænger, fortrængte, har fortrængt)} & 排挤，取代，驱离 \\
动词 [har] & at presse & \textit{(presser, pressede, har presset)} & 施压，推高，挤迫 \\
动词 [har] & at udbygge & \textit{(udbygger, udbyggede, har udbygget)} & 扩建，拓展完善 \\
动词 [har] & at bremse & \textit{(bremser, bremsede, har bremset)} & 抑制，放缓，刹住 \\
\midrule
形容词 & eftertragtet & \textit{(eftertragtet, eftertragtede)} & 抢手的，备受追捧的 \\
形容词 & almen & \textit{(alment, almene)} & 公益性的，非营利福利性质的 \\
副词短语 & paradoksalt nok & \textit{(不变，句首评价副词)} & 令人悖论的是，具有讽刺意味的是 \\
\end{glosertable}

% ---------------------------------------------------------------------
\ncdsubsec{四、Grammatik \& Stilistik（动态同步、比例递进与悖论句法精讲）}

\begin{grammatikbox}{语法点 1：动态同步、比例递进与悖论修辞总览矩阵}
在 C1 级别的社会评论与城市分析文章中，静态列举原因往往不够立体；用好下面五组句式，才能让论述呈现“随时间同步演化、相互加速”的动态张力：

\vspace{0.3em}
\renewcommand{\arraystretch}{1.24}
{\small
\begin{tabularx}{\linewidth}{@{} >{\bfseries\raggedright\arraybackslash}p{2.9cm} >{\raggedright\arraybackslash}p{4.6cm} >{\raggedright\arraybackslash}X @{}}
\toprule
\textbf{修辞与逻辑功能} & \textbf{句型结构公式} & \textbf{高阶社会评论例句} \\
\midrule
\dabold{1. 悖论揭示} &
\textit{\dabold{Paradoksalt nok} + V2 倒装} &
\textit{\dabold{Paradoksalt nok} fører byfornyelse ofte til mindre social mangfoldighed.} \\
\midrule
\da{2. 动态同步推进} &
\textit{\da{I takt med at} [主+动]..., [动+主]...}\newline
（或 \textit{i takt med} + 名词） &
\textit{\da{I takt med at} byen bliver grønnere, stiger boligpriserne hurtigt.} \\
\midrule
\da{3. 比例同步递进} &
\textit{\da{Jo} + 比较级 [从句语序], \da{desto} + 比较级 [V2 倒装]} &
\textit{\da{Jo mere} attraktiv bydelen bliver, \da{desto} højere presses huslejen op.} \\
\midrule
\dabold{4. 正反馈恶性循环} &
\textit{sætte \dabold{en selvforstærkende effekt} i gang / bryde...} &
\textit{Høje priser sætter \dabold{en selvforstærkende effekt} i gang på boligmarkedet.} \\
\midrule
\dabold{5. 代价与利益取舍} &
\textit{\dabold{på bekostning af} + 名词短语} &
\textit{Luksusbyggeri sker ofte \dabold{på bekostning af} billige almene boliger.} \\
\bottomrule
\end{tabularx}
}
\end{grammatikbox}

\begin{grammatikbox}{语法点 2：\textit{Jo ..., desto ...} 语序铁律与三大考场陷阱}
\begin{enumerate}
  \item \textbf{\da{Jo ..., desto ...} 的核心秘诀：“前从句不倒装，后主句必倒装”！}
    \begin{itemize}
      \item 前半句由 \textit{jo} 引导，属于\textbf{从属从句}，内部保持“主语 +（\textit{ikke}）+ 动词”正常语序；当“\textit{jo + 比较级 + 名词}”本身充当从句主语时，后面必须插入形式主语占位词 \dabold{der}：\\
      \ding{51}\ \textit{Jo flere familier \dabold{der} flytter ud, \da{desto} længere bliver pendlertiden.}
      \item 后半句由 \textit{desto} 引导，属于\textbf{主句}，而“\textit{desto + 比较级（+ 被修饰名词）}”整个词组占领句首第一顺位，因此后面\textbf{必须立刻接变位动词}（V2 倒装）：\\
      \ding{51}\ \textit{Jo dyrere boligerne bliver, \da{desto sværere får} unge lønmodtagere ved at blive.}（绝不可写成 \ding{55}\ \textit{*...desto unge får sværere ved...}）。
    \end{itemize}
  \item \textbf{\textit{Paradoksalt nok} 与 \textit{I takt med at...} 不加逗号，且强触发 V2 倒装}：\\
    \ding{51}\ \textit{Paradoksalt nok \dabold{har} sygeplejersker ikke råd til at bo tæt på hospitalet.}（英语常写 \textit{Paradoxically, ...}，但丹麦语\textbf{绝不能在句首副词后打逗号}！）。
  \item \textbf{\textit{på bekostning af} 的固定形态}：表示“以……为代价”，永远用单数无冠词形式 \textit{på bekostning af} + 名词短语；不要误加定冠词写成 \ding{55}\ \textit{*på bekostningen af} 或改用介词 \ding{55}\ \textit{*på bekostning for}。
\end{enumerate}
\end{grammatikbox}

% ---------------------------------------------------------------------
\ncdsubsec{五、Mønsterøvelser（悖论与动态结构操练）}

\textbf{练习 A：动态结构与核心搭配精准填空（每空填入最准确的 1 个丹麦语词）}
\begin{enumerate}
  \item I \underline{\hspace{2.2cm}} med at storbyen tiltrækker flere studerende, vokser presset på boligmarkedet.
  \item \underline{\hspace{2.2cm}} nok fører grøn byfornyelse ofte til, at almindelige familier presses ud af kvarteret.
  \item Jo højere kvadratmeterprisen stiger, \underline{\hspace{2.2cm}} sværere bliver det for unge at købe lejlighed.
  \item Jo flere pendlere \underline{\hspace{2.2cm}} bosætter sig i oplandet, desto større bliver trængslen på motorvejen.
  \item Eksklusive byggerier prioriteres desværre ofte på \underline{\hspace{2.2cm}} af billige almene boliger.
  \item Når velstillede købere dominerer et område, sætter det en \underline{\hspace{2.2cm}} effekt i gang.
  \item Kommunen ønsker at \underline{\hspace{2.2cm}} segregeringen gennem faste kvoter for blandede ejerformer.
  \item Hurtigere kommunal \underline{\hspace{2.2cm}} kan gøre det lettere at bygge flere betalelige boliger i tide.
\end{enumerate}

\textbf{练习 B：动态比例句语序纠错与句型升级（修正语病或按要求重构）}
\begin{enumerate}
  \item 【纠错】修正英语逗号与 V2 倒装错误：\textit{Paradoksalt nok, manglen på billige boliger bliver værre i grønne bydele.} $\rightarrow$
  \item 【纠错】修正 \textit{jo...desto...} 的主从句语序与占位词：\textit{Jo flere dyre lejligheder bygges i centrum, desto færre lærere kan bo nær skolen.} $\rightarrow$
  \item 【纠错】修正介词短语错用：\textit{Byen vokser hurtigt på bekostningen for den sociale mangfoldighed.} $\rightarrow$
  \item 【升级】用 \textit{I takt med at...} 合并两句：\textit{Byerne investerer i nye metrostationer. Samtidig stiger huslejen hurtigt.} $\rightarrow$
  \item 【升级】用 \textit{Jo ..., desto ...} 重写因果递进句：\textit{Hvis borgere flytter længere væk fra jobbet, bruger de mere tid på daglig transport.} $\rightarrow$
  \item 【升级】用 \textit{sætte en selvforstærkende effekt i gang} 浓缩说明：\textit{Når et kvarter bliver populært, stiger priserne, og det gør kvarteret endnu mere eksklusivt.} $\rightarrow$
\end{enumerate}

\textbf{练习 C：中译丹（活用悖论、同步演化与比例递进句式）}
\begin{enumerate}
  \item 随着大城市变得越来越宜居，普通家庭往往负担不起市中心的房租。
  \item 令人悖论的是，成功的老旧街区改造有时会削弱城市的社会多样性。
  \item 市中心建得越多高价公寓，外围市镇承受的通勤压力就越大。
  \item 奢华住房项目常常是以牺牲平价福利住房为代价推进的。
  \item 当富裕买家接管一个街区时，就会启动一种自我强化效应。
  \item 城市规划必须将经济增长与公共交通结合起来，才能遏制社会隔离。
\end{enumerate}

% ---------------------------------------------------------------------
\ncdsubsec{六、PD3 / Studieprøven 实战微任务与参考答案}

\begin{emailbox}{Studieprøven / PD3 Debatindlæg 实战微任务（Samfundsanalyse · mindst 220--250 ord）}
\textbf{Situation}: I en avisdebat om boligpolitikken i de danske storbyer skriver en økonom, at markedet alene bør bestemme, hvad der bygges, fordi flere nye boliger automatisk vil sænke priserne for alle.\\
\textbf{Opgave}: Skriv et analytisk debatindlæg på \textbf{mindst 220--250 ord}, hvor du diskuterer urbaniseringens pris i storbyerne. Du skal:
\begin{enumerate}
  \item Analysere paradokset omkring stigende attraktivitet og sociale udfordringer i centrum.
  \item Forklare den selvforstærkende effekt og konsekvenserne for pendlerne i oplandskommunerne.
  \item Tage stilling til økonomens påstand og foreslå konkrete politiske tiltag, der kan sikre blandede bydele (brug gerne \textit{paradoksalt nok, i takt med at, jo ..., desto ...} og \textit{på bekostning af}).
\end{enumerate}
\end{emailbox}

\begin{answerbox}
\textbf{Før du læser 参考回答}
\begin{enumerate}
  \item \textit{Paradokset består i, at netop de vellykkede grønne og kulturelle investeringer får grundværdierne til at stige så kraftigt, at pædagoger og sygeplejersker prissættes ud af byen. Boligmarkedet skaber en selvforstærkende effekt ved, at eksklusive kvarterer tiltrækker flere velstillede købere og investorer, hvilket presser kvadratmeterpriserne endnu højere op.}
  \item \textit{Henrik frygter, at strenge kvoter på 30 procent almene boliger vil bremse privat nybyggeri og skabe boligmangel, mens Sofie mener, at dyrt byggeri alene ikke hjælper almindelige lønmodtagere. Til sidst mødes de om hurtigere sagsbehandling, stationsnær udbygning og blandede ejerformer.}
\end{enumerate}

\textbf{Tekst B 阅读核对答案}
\begin{enumerate}
  \item \da{Rigtigt}（他明确担忧强制配额会吓退私人投资者，导致建房总数减少）
  \item \da{Fordi huslejen ligger over 16.000 kroner, hvilket almindelige lønmodtagere ikke kan betale.}
  \item \da{En hurtigere kommunal sagsbehandling kombineret med udbygning omkring nye S-togsstationer i forstæderne.}
\end{enumerate}

\textbf{练习 A 答案}
\begin{enumerate}
  \item \da{takt}（固定动态同步连接短语：\textit{i takt med at...} 随着……同步发展）
  \item \da{Paradoksalt}（句首评价副词短语：\textit{Paradoksalt nok} 令人悖论的是）
  \item \da{desto}（比例比较结构配对：\textit{jo ..., desto ...}）
  \item \da{der}（当 \textit{jo flere + 名词} 在从句中作主语时，必须接占位关系词 \textit{der}）
  \item \da{bekostning}（固定搭配：\textit{på bekostning af} 以牺牲……为代价）
  \item \da{selvforstærkende}（固定学术搭配：\textit{en selvforstærkende effekt} 自我强化效应）
  \item \da{bremse} / \da{modvirke}（动词不定式：遏制/抑制社会隔离）
  \item \da{sagsbehandling}（名词：市政行政审批流程）
\end{enumerate}

\textbf{练习 B 答案（参考）}
\begin{enumerate}
  \item \textit{Paradoksalt nok bliver manglen på billige boliger værre i grønne bydele.}（删去英语式逗号，并保证变位动词 \textit{bliver} 紧跟在 \textit{Paradoksalt nok} 之后）
  \item \textit{Jo flere dyre lejligheder der bygges i centrum, desto færre lærere kan bo nær skolen.}（前句加入主语占位词 \textit{der}；后句将 \textit{desto færre lærere} 视为句首成分，紧接倒装动词 \textit{kan}）
  \item \textit{Byen vokser hurtigt på bekostning af den sociale mangfoldighed.}（固定为单数无冠词 \textit{på bekostning af}）
  \item \textit{I takt med at byerne investerer i nye metrostationer, stiger huslejen hurtigt.}（以 \textit{I takt med at...} 体现发展同步性，并注意主句 V2 倒装 \textit{stiger huslejen}）
  \item \textit{Jo længere væk fra jobbet borgerne flytter, desto mere tid bruger de på daglig transport.}（严格遵守“前从句不倒装、后主句 \textit{bruger de} 倒装”）
  \item \textit{Når et kvarter bliver populært, sætter prisstigningerne en selvforstærkende effekt i gang.}（用核心词组精炼概括正反馈螺旋）
\end{enumerate}

\textbf{练习 C 答案}
\begin{enumerate}
  \item \textit{I takt med at storbyerne bliver mere attraktive, har almindelige familier ikke råd til huslejen i centrum.}
  \item \textit{Paradoksalt nok kan en vellykket byfornyelse svække byens sociale mangfoldighed.}
  \item \textit{Jo flere dyre lejligheder der bygges i centrum, desto større bliver pendlerpresset i oplandskommunerne.}
  \item \textit{Luksusprojekter gennemføres ofte på bekostning af billige almene boliger.}
  \item \textit{Når velstillede købere overtager et kvarter, sættes en selvforstærkende effekt i gang.}
  \item \textit{Byplanlægningen skal kombinere vækst med offentlig transport for at bremse segregeringen.}
\end{enumerate}

\textbf{PD3 / Studieprøven 高分范文（约 240 词，不含标题）}\\
\textit{\textbf{Markedet alene skaber ikke en by for alle}}

\textit{Det lyder umiddelbart besnærende, når en økonom hævder, at et helt frit boligmarked automatisk vil dæmpe prispresset gennem øget nybyggeri. Erfaringerne fra København og Aarhus tegner imidlertid et langt mere nuanceret og bekymrende billede. \da{I takt med at} storbyerne gennem de seneste årtier er blevet grønnere, tryggere og kulturelt mere levende, er kvadratmeterpriserne steget langt hurtigere end almindelige lønninger.}

\textit{\dabold{Paradoksalt nok} risikerer selve byens succes dermed at undergrave dens sociale mangfoldighed. Når private investorer udelukkende bygger efter det højeste afkast, opføres der primært dyre ejer- og lejeboliger \dabold{på bekostning af} billige almene boliger. Det sætter \dabold{en selvforstærkende effekt} i gang på boligmarkedet: \da{Jo mere} eksklusivt et kvarter bliver, \da{desto} færre pædagoger, sygeplejersker og politibetjente har råd til at blive boende tæt på deres arbejdsplads og de borgere, de servicerer.}

\textit{Resultatet er en stigende pendlerulighed mellem bykernen og omegnen. \da{Jo længere} væk fra arbejdspladsen almindelige lønmodtagere fortrænges, \da{desto} flere timer og ressourcer mister familierne på daglig transport i myldretiden. Et rent frit marked løser altså ikke denne skæve udvikling, fordi efterspørgslen efter attraktive adresser bliver ved med at presse huslejen opad.}

\textit{Derfor har kommunerne brug for en aktiv og langsigtet planpolitik. Ved at kombinere hurtigere sagsbehandling med klare krav til blandede ejerformer, faste kvoter for almene boliger og stærke S-togsforbindelser til oplandskommunerne kan byerne både rumme fortsat vækst og sikre, at almindelige borgere også i fremtiden har råd til at bo i byen.}

\small\textbf{名师点评}：这篇约 240 词的评论文章精准回应了经济学家观点并剖析了城市化代价。全篇高频且无误地驾驭了本课五大支柱：\textit{I takt med at...}（同步推移）、\textit{Paradoksalt nok}（首句倒装揭示悖论）、\textit{på bekostning af}（利益权衡）、\textit{en selvforstærkende effekt}（机制提炼）以及两处高难度的 \textit{Jo ..., desto ...} 句式，兼具社会学洞察与 C1 句法掌控力。
\end{answerbox}
